% SGA 3, Expose VI_B, Loop 1 English reconstruction.
% Scope: Proposition 3.7 and Corollary 3.8 in full.
% Authority: Polo--Gille Exp6B-13oct24.pdf, component-local pages 19--20
% (combined-reader pages 360--361; original-margin printed pages 343--345).
% OCR and the Reinhold English Markdown are comparison witnesses only.

\medskip
\phantomsection\label{prop:sga3-VIB-3.7}%
\noindent\textbf{命题 3.7.---}%
\footnote{\label{note:sga3-VIB-24}编者注: 回顾以下定义和结果
(参见 EGA~0\(_{\mathrm{III}}\), \S9.1 及 EGA~IV\(_1\), \S\S1.8 和 1.9).
设 \(X\) 为拓扑空间.
\textup{(i)} \(X\) 的开子集 \(U\) 称为\emph{逆紧的 (retrocompact)},
是指包含映射 \(U\hookrightarrow X\) 拟紧,
即对每个拟紧开子集 \(V\subset X\), \(U\cap V\) 都拟紧.
\textup{\SGARefLink{prop:sga3-VIB-3.7:item:ii}{(ii)}}
\(X\) 的子集 \(C\) 称为\emph{可构造的}, 是指它是有限个形如
\(U-V\) 的子集的并, 其中 \(U\) 和 \(V\) 是 \(X\) 的逆紧开子集.
\textup{\SGARefLink{prop:sga3-VIB-3.7:item:iii}{(iii)}}
\(X\) 的子集 \(L\) 称为\emph{局部可构造的}, 是指对每个 \(x\in X\),
存在 \(x\) 在 \(X\) 中的开邻域 \(U\), 使 \(L\cap U\) 在 \(U\) 中可构造.
(若 \(X\) 拟紧且拟分离, 则 \(L\) 可构造.)
\textup{(iv)} \(X\) 的子集 \(I\) 称为\emph{归纳可构造的 (ind-constructible)},
是指对每个 \(x\in X\), 存在 \(x\) 在 \(X\) 中的开邻域 \(U\),
使 \(I\cap U\) 是 \(U\) 的局部可构造子集的并.
现设 \(f:X\to Y\) 为概形态射. 由 Chevalley 可构造性定理
(参见 EGA~IV\(_1\), 1.8.4 和 1.9.5(viii)),
若 \(f\) 为有限呈示 (分别地, 局部有限呈示),
则每个局部可构造子集在 \(f\) 下的像都局部可构造
(分别地, 归纳可构造).}
\emph{设 \(G\) 为局部有限呈示的 \(S\)-群. 则:}
\begin{enumerate}
\item[\textup{\SGAInlineRefTarget{prop:sga3-VIB-3.7:item:i}(i)}]
  \emph{\(\underline{G^0}\) 在 \(G\) 中归纳可构造.}
\item[\textup{\SGAInlineRefTarget{prop:sga3-VIB-3.7:item:ii}(ii)}]
  \emph{若此外 \(G\) 在 \(S\) 上拟分离, 且 \(S\) 拟紧且拟分离,
  则 \(\underline{G^0}\) 可构造.}
\item[\textup{\SGAInlineRefTarget{prop:sga3-VIB-3.7:item:iii}(iii)}]
  \emph{因此, 若 \(G\) 在 \(S\) 上拟分离,
  则 \(\underline{G^0}\) 局部可构造.}
\end{enumerate}

先证明第一个断言. 由于
\[
  \pi:G\longrightarrow S
\]
局部有限呈示, 给定 \(s\in S\), 存在 \(G\) 中包含 \(\epsilon(s)\)
的开子集 \(U\), 以及 \(S\) 中包含 \(s\) 的开子集 \(V\), 使
\[
  \pi(U)\subset V
\]
且诱导态射
\[
  \pi':U\longrightarrow V
\]
为有限呈示. 则
\[
  T=\epsilon^{-1}(U)
\]
是 \(S\) 的开子集. 令
\[
  W=(\pi')^{-1}(T),\qquad
  \pi''=\pi'|_W.
\]
于是 \(\pi'':W\to T\) 为有限呈示,
且具有由 \(\epsilon\) 诱导的态射
\[
  \epsilon'':T\longrightarrow W
\]
作为截面.

对每个 \(t\in T\), 由于 \(\underline{G_t^0}\) 不可约
(VI\(_A\), \hyperref[prop:sga3-VIA-2.4]{2.4}),
\(W\cap\underline{G_t^0}\) 在 \(\underline{G_t^0}\) 中稠密,
因而不可约, 故连通.
因此它是 \((\pi'')^{-1}(t)\) 中包含 \(\epsilon''(t)\) 的连通分支.
由 EGA~IV\(_3\), 9.7.12 可知,
当 \(t\in T\) 时各 \(W\cap\underline{G_t^0}\) 的并 \(W^0\)
在 \(W\) 中局部可构造.
另一方面, VI\(_A\), \SGARefLink{sga3:vi:proposition:0-5}{0.5} 给出
\[
  \underline{G^0}\cap\pi^{-1}(T)=W^0\cdot W^0;
\]
即 \(\underline{G^0}\cap\pi^{-1}(T)\) 是
\[
  W^0\times_T W^0
\]
在 \(\mu\) 诱导的态射
\[
  \mu'':W\times_T W\longrightarrow\pi^{-1}(T)
\]
下的像.%
\footnote{\label{note:sga3-VIB-25}编者注: 以下展开了原文.}
由于 \(W\times_T W\) (分别地, \(\pi^{-1}(T)\))
在 \(T\) 上为有限呈示 (分别地, 局部有限呈示),
由 EGA~IV\(_1\), 1.4.3 和 1.2.2,
态射 \(\mu''\) 局部有限呈示且拟分离.
若此外 \(\pi^{-1}(T)\) 在 \(T\) 上拟分离,
则 \(\mu''\) 拟紧 (上述文献, 1.2.4), 因而为有限呈示.
由于 \(W^0\times_T W^0\) 在 \(W\times_T W\) 中局部可构造
(因为 \(W^0\) 在 \(W\) 中局部可构造),
Chevalley 可构造性定理 (上述文献, 1.8.4 和 1.9.5(viii)) 表明,
\(\underline{G^0}\cap\pi^{-1}(T)\) 在 \(\pi^{-1}(T)\) 中归纳可构造;
且若 \(G\), 从而 \(\pi^{-1}(T)\), 在 \(T\) 上拟分离,
则它在 \(\pi^{-1}(T)\) 中局部可构造.
这证明了断言 \textup{\SGARefLink{prop:sga3-VIB-3.7:item:i}{(i)}} 和
\textup{\SGARefLink{prop:sga3-VIB-3.7:item:iii}{(iii)}}.

现假设 \(G\) 在 \(S\) 上拟分离, 且 \(S\) 拟紧且拟分离.
由 \hyperref[prop:sga3-VIB-3.5]{3.5},
存在 \(G\) 的拟紧开子集 \(U\), 包含 \(\underline{G^0}\).
由于 \(G\) 在 \(S\) 上拟分离, \(G\) 拟分离;
因而开子集 \(U\) 逆紧 (EGA~IV\(_1\), 1.2.7),
只须说明 \(\underline{G^0}\) 在 \(U\) 中可构造
(EGA~0\(_{\mathrm{III}}\), 9.1.8).

此外, \(U\) 拟紧, 因而在 \(S\) 上拟紧 (EGA~IV\(_1\), 1.2.4),
且它在 \(S\) 上拟分离.
故 \(\pi\) 在 \(U\) 上的限制为有限呈示.
由 EGA~IV\(_3\), 9.7.12, \(\underline{G^0}\) 在 \(U\) 中局部可构造,
因而在 \(U\) 中可构造, 因为 \(U\) 拟紧且拟分离
(EGA~IV\(_1\), 1.8.1).
这证明了 \textup{(ii)}, 并且可知对 \(S\) 的每个拟紧且拟分离开子集
\(T\) (例如 \(S\) 的每个仿射开子集),
\[
  \underline{G^0}\cap\pi^{-1}(T)
\]
都可构造.

\medskip
\phantomsection\label{cor:sga3-VIB-3.8}%
\noindent\textbf{推论 3.8.---}
\emph{设 \(S_0\) 为拟紧且拟分离的概形,
\(I\) 为递增滤过的预序集,
\((\mathcal A_i)_{i\in I}\) 为交换的拟凝聚 \(\mathcal O_{S_0}\)-代数的正向系,}
\[
  \mathcal A=\varinjlim_i\mathcal A_i,\qquad
  S_i=\operatorname{Spec}\mathcal A_i\quad(i\in I),\qquad
  S=\operatorname{Spec}\mathcal A
\]
\emph{\textup{(参见 EGA~II, 1.3.1)}.
设 \(G\) 为局部有限呈示的 \(S_0\)-群概形. 则典范映射}
\[
  \varinjlim_i G^0(S_i)\longrightarrow G^0(S)
\]
\emph{是双射.}

由于 \(G\) 在 \(S_0\) 上局部有限呈示,%
\begingroup
\renewcommand{\thefootnote}{\fnsymbol{footnote}}%
\footnote[1]{\label{note:sga3-VIB-3.8-base}译者注:
法文底本在此印作``在 \(S\) 上''.
概形 \(G\) 定义在 \(S_0\) 上, 而 \(S\) 及各 \(S_i\)
是正向系中的测试概形; 因此所需基底为 \(S_0\).}%
\endgroup{}
典范映射
\[
  \varinjlim_i G(S_i)\longrightarrow G(S)
\]
由 EGA~IV\(_2\), 8.14.2(c) 为双射.
立即可知典范映射
\[
  \varinjlim_i G^0(S_i)\longrightarrow G^0(S)
\]
为单射. 证明其满射. 设
\[
  g\in G^0(S)\subset G(S).
\]
存在 \(i\in I\), 使 \(g\) 借助元素 \(g_i\in G(S_i)\)
经过 \(S\to S_i\) 分解. 由假设,
\[
  g^{-1}(\underline{G^0})=S.
\]
但由 \hyperref[prop:sga3-VIB-3.7]{3.7},
\(\underline{G^0}\) 在 \(G\) 中归纳可构造, 因而
\[
  g_i^{-1}(\underline{G^0})
\]
在 \(S_i\) 中归纳可构造.
由 EGA~IV\(_2\), 8.3.4, 存在指标 \(j\geq i\), 使
\[
  g_j^{-1}(\underline{G^0})=S_j,
\]
其中 \(g_j\) 是 \(g_i\) 经基变换 \(S_j\to S_i\) 诱导的映射.
因此 \(g_j\in G^0(S_j)\), 从而 \(g\) 来自
\[
  \varinjlim_i G^0(S_i).
\]
的某个元素.
% Translation statistics: source_lines=201; translated_lines=189; inline_math=148; display_math=19; pending_TODO=0.
% Review: complete component translation; math fragments, reference keys, environments and note counts checked against source; formula prose translated; compilation pending parent integration.
